% ============================================================
%  SECTION 4 — Shor's Algorithm
% ============================================================
\section{Shor's Algorithm: A Step-by-Step Walkthrough}

Let's first state clearly the goal of the algorithm:
\textbf{Given a semiprime number $N$, find the two prime numbers $P_1$ and
$P_2$ such that: $P_1 \times P_2 = N$.}

The classical way to find these numbers is just by guessing randomly, also
known as the brute force method. The essence or the main idea of Shor's
algorithm is to translate the problem of integer factorization, into a problem
of period finding, and to show that given a random guess $g$, this bad guess
can be systematically transformed into a better one.

\subsection{Initialization}

We start by guessing randomly a number $g$ such that $1 < g < N$. This leads
to two cases:

\textbf{Case 1:} We got lucky. $g$ is a factor of $N$, or shares common
factors with $N$. In this case, we are done thanks to Euclid's algorithm.

\begin{smjdefinition}
\textbf{Euclid's Algorithm} is a classical method for efficiently computing
the \textbf{greatest common divisor ($GCD$)} of two integers. It works by
repeatedly replacing the larger number with its remainder when divided by the
smaller, until the remainder is zero. The last non-zero remainder is the $GCD$.
\end{smjdefinition}

\begin{smjnote}
Using that, to find a factor of $N$, we don't need then to guess a factor,
just to find a number that shares factors with $N$.
\end{smjnote}

After finding one of the prime numbers, let's assume $P_1$, compute $N/P_1$
and this will give you $P_2$.

\textbf{Case 2:} Most probably, we won't get lucky.
And this will make us start our process of making our bad guess a better one.
To achieve that, we will borrow a quite helpful theorem.

\begin{smjtheorem}
For any two whole numbers $A$ and $B$ that are coprime, there exists a power
$r$ and a multiple $m$ such that:
$$A^r = m \cdot B + 1$$
Equivalently:
$$A^r \equiv 1\pmod B$$
\end{smjtheorem}

A good example to visualize it will be $A = 2$ and $P = 5$, we find that
$r = 4$:
$$
2^1 \equiv 2 \pmod{5} \qquad
2^2 \equiv 4 \pmod{5} \qquad
2^3 \equiv 3 \pmod{5} \qquad
2^4 \equiv 1 \pmod{5}
$$

In our case, our $A$ is our bad guess $g$, and $B$ is our semiprime number
$N$ that we want to factor, and ultimately, we know the existence of our $r$.

\subsection{Making the Bad Guess a Better Guess}

Let's assume that we found $r$, such that $g^r \equiv 1 \pmod{N}$. Bring $1$
to the left and you get:
$$
g^r - 1 \equiv 0\pmod{N}
$$

And using our famous identity: $(a + b)(a - b) = a^2 - b^2$:

Replace $a$ by $g$, and $b$ by $1$, and going from the right side of the
equality to the left, we get:
$$
\underbrace{(g^{r/2} - 1)}_{\text{$P_1'$}} \cdot
\underbrace{(g^{r/2} + 1)}_{\text{$P_2'$}} \equiv 0\pmod{N}
$$

Notice that $r$ has to be even, otherwise, $r/2$ won't be a whole number. If
we found an $r$ odd, we repeat our algorithm from the beginning.

What we found now is actually really close to what we want, instead of finding
the two numbers $P_1$ and $P_2$ such that $P_1 \times P_2 = N$, we found that
$N$ divides the product of two other numbers $P_1'$ and $P_2'$. Which can be
written as:
$$
P_1' \times P_2' = mN \quad , m \in \mathbb{N}
$$

Even more precisely, we found two numbers $P_1'$ and $P_2'$ that share factors
with $N$. Thus:
$$
\gcd(P_1', N) = P_1 \text{ or } P_2.
\text{ (We can replace } P_1' \text{ by } P_2', \text{ the outcome stays the same).}
$$
And remember that thanks to our beautiful Euclid's algorithm, we are able to
compute $\gcd(P_1', N) = P_1$ easily. Compute then $\dfrac{N}{P_1} = P_2$,
(and vice versa if we got that $\gcd(P_1', N) = P_2$). And we have
successfully found our $P_1$ and $P_2$.

\textit{End of our algorithm.}

Or maybe not?

You may feel that this is too easy, in the sense where, using this algorithm
we could directly break the RSA, why are we using it? Remember also that I
gave you an introduction about quantum computing, but I didn't talk at all
about it at any point.

Let's recall the algorithm's main steps.

\begin{tcolorbox}[colframe=smjcrimson, colback=smjlightcrimson,
                  sharp corners, boxrule=0.5pt, breakable]
1 -- Make a guess, $g < N$ that shares no factors with $N$.\\[0.5em]
2 -- Find $r$ such that $g^r = mN + 1$.\\[0.5em]
3 -- If $r$ is even, calculate $(g^{r/2} + 1)$ and $(g^{r/2} - 1)$.
     If $r$ is odd, go back to step 1.\\[0.5em]
4 -- Use Euclid's algorithm to find the greatest common divisor.
\end{tcolorbox}

Where is the \textbf{quantum} here?

\subsection{The Quantum Subroutine of Shor's Algorithm}

Notice that we didn't discuss before how we are finding $r$, which is in fact
the hardest (probably due to the quantum concepts) part of the whole algorithm.

\begin{smjnote}
Through this part, the quantum part of the explanation (which is still the
main thing) is going to be a bit simplified, essentially when we come to some
deep notions such as Quantum Fourier Transform, etc. Such quantum concepts as
said before need even their own whole article. And as again mentioned before,
one of the primary goals of the article is to keep the reader on board.
There will be some analogies to the classical architecture of a computer, for
the sake of your understanding.
\end{smjnote}

\subsubsection{Initialization}

To achieve this quantum subroutine, we will need two quantum registers, as
normal registers that we have in normal computers, these ones are just going
to contain qubits instead of normal bits.

First register: is going to be $2n$ qubits, initialized as:
\[
\ket{0}^{\otimes 2n}
\]
The number $n$ just determines how accurate we want our calculations to be,
however, the 2nd register size depends on the first one.

Second register is going to be $n$ qubits, initialized to: $\ket{1}$

Note that $\otimes$ in context of quantum computing doesn't mean the XOR gate
as a lot of people would think. It is a tensor product. Just consider it as a
normal multiplication but between the quantum states.

Note also that for the first register, we mean that all the bits are set to
$0$, meaning that the state of that register is: $\ket{000\ldots000}$.

But for the second register, we mean that the value of that register is one,
meaning that the state of that register is: $\ket{000\ldots001}$. It is just
a story of notations.

\subsubsection{Hadamard Transform}

Let's first define how the Hadamard Transform works. The Hadamard gate is
formally defined as this:
$$
H = \frac{1}{\sqrt{2}}
\begin{bmatrix}
1 & 1 \\
1 & -1
\end{bmatrix}
$$

What does it do? Applying the Hadamard Transform to a $\ket{1}$ or $\ket{0}$
gives this:
\begin{align*}
H\ket{0} = \frac{1}{\sqrt{2}}(\ket{0} + \ket{1})\\
H\ket{1} = \frac{1}{\sqrt{2}}(\ket{0} - \ket{1})
\end{align*}

During this part, we will be interested essentially in its application on the
$0$ state, since our first register is just full of $0$'s.

Notice that on $0$, it created exactly a state such that the probability of
getting one of the states $x$ such that $x \in \{0, \ldots, 2^n - 1\}$
(converting the decimal to binary) is exactly equal.

Let's try to apply it to $0$ states with $2$ qubits.
$$
H^{\otimes 2} \ket{00} = \frac{1}{2}(\ket{00} + \ket{01} + \ket{10} + \ket{11})
$$

And on $3$ qubits?
$$
H^{\otimes 3} \ket{000} = \frac{1}{\sqrt{8}} (\ket{000} + \ket{001} + \ket{010}
+ \ket{011} + \ket{100} + \ket{101} + \ket{110} + \ket{111})
$$
Did you notice the pattern? If not let me show it to you:

Whenever we apply the $H$ gate to a register of size $n$, such that all its
qubits are $0$, we get a state such that the probability of measuring a state
in $\{0, 1, 2, \ldots, 2^n - 1\}$ is exactly equal for all of them.

You may wonder now why we are doing that, we will see why soon.

To end this part, we will formally say that the state of our first register
now is:
$$
\frac{1}{\sqrt{2^n}} \sum_{x=0}^{2^n - 1} \ket{x}
$$

\subsubsection{$U$}

We will apply another transformation to the full two-register system. Let's
call this transformation $U$. Formally, $U$ acts on a pair of states
$\ket{x} \otimes \ket{y}$ as follows:
$$
U(\ket{x} \otimes \ket{y}) = \ket{x} \otimes \ket{y \cdot g^x \bmod N}
$$
With $g$ the random guess that we picked before. In other words, $U$ reads the
value $x$ from the first register and uses it to update the second register:
it multiplies the second register's value by $g^x$, modulo $N$.

How are we going to use $U$? First let's look at our state now.

The quantum state of our full system (including the $1^\text{st}$ and the
$2^\text{nd}$ register) is the tensor product $\otimes$ of both first and
second register. As I said $\otimes$ is just a normal multiplication of the
states. Our first register is now a sum of equal states, and our second
register is still just in a $1$ state. ``Expanding'' the expression:

\[
\left( \frac{1}{\sqrt{2^n}} \sum_{x=0}^{2^n - 1} \ket{x} \right) \otimes \ket{1}
\]

gives:

\[
\frac{1}{\sqrt{2^n}} \sum_{x=0}^{2^n - 1} \left( \ket{x} \otimes \ket{1} \right)
\]

Don't overthink about it, it is the exact same way as you would expand
$(a + b + c)d = ad + bd + cd$.
For each $x$ (the values in each state from $0$ to $2^n - 1$), we apply the
transformation $x$ times on the qubit at state $1$ associated to it. It is a
bit hard to understand it like that, we can use a visualization to understand
it more:

For each pair $\ket{x} \otimes \ket{1}$, we apply $U$ getting:
$$
\ket{x} \otimes \ket{1} \longrightarrow \ket{x} \otimes U^x \ket{1}
= \ket{x} \otimes \ket{g^x \bmod N}
$$

As an example, applying $U$ to a state $\ket{101} \otimes \ket{1}$
($101$ read as $4 + 0 + 1 = 5$) gives:
$$
\ket{101} \otimes U^5 \ket{1} = \ket{101} \otimes \ket{g^5 \bmod N}
$$
And again, $g$ is our random guess that we picked before.

\begin{smjnote}
$\ket{x}\otimes\ket{y}$ is the same as $\ket{x}\ket{y}$
\end{smjnote}

Notice that when applying the transformation, we will get a certain periodicity
in our quantum system due to modulo properties. To visualize it even better,
let us consider that our $x$ is going from $0$ to $7$ (using thus $3$ qubits),
that our $N$ was $26$, and our random guess was $3$, after applying everything
we had before, we will get:
$$
\frac{1}{\sqrt{8}} (\ket{000} \ket{3^0 \bmod 26} + \ket{001} \ket{3^1 \bmod 26}
+ \ket{010} \ket{3^2 \bmod 26} +
$$
$$
\ket{011} \ket{3^3 \bmod 26} + \ket{100} \ket{3^4 \bmod 26}
+ \ket{101} \ket{3^5 \bmod 26} +
$$
$$
\ket{110} \ket{3^6 \bmod 26} + \ket{111} \ket{3^7 \bmod 26})
$$

Which exactly gives:
$$
\frac{1}{\sqrt{8}} (\ket{000} \ket{1 \bmod 26} + \ket{001} \ket{3 \bmod 26}
+ \ket{010} \ket{9 \bmod 26} +
$$
$$
\ket{011} \ket{1 \bmod 26} + \ket{100} \ket{3 \bmod 26}
+ \ket{101} \ket{9 \bmod 26} +
$$
$$
\ket{110} \ket{1 \bmod 26} + \ket{111} \ket{3 \bmod 26})
$$

Notice also that the period is $3$, which is exactly the number $r$ that we
are searching for, $3^3 \equiv 1 \pmod{26}$.

Thus, if we are able to extract the periodicity from the state, we would be
done.

\subsubsection{The Periodicity, a.k.a.\ $r$}

This is actually a much harder task than expected. When you see the
representation of the state in front of you it seems obvious, however, as we
said before, if we just measure now our state it is just going to give us a
random state since they all have equal probability. Which is not what we want.
In fact, we are not able to extract the periodicity. However, we have something
called the QFT (for \textbf{Q}uantum \textbf{F}ourier \textbf{T}ransform).

The QFT is the quantum analogue of the classical Discrete Fourier Transform
(DFT), which is a well-known tool in signal processing for decomposing a
signal into its frequency components. The DFT of a sequence
$x_0, x_1, \ldots, x_{N-1}$ is defined as:
$$
X_k = \sum_{n=0}^{N-1} x_n \cdot e^{\frac{2\pi i k n}{N}}
$$
The QFT is defined as:
$$
\text{QFT} \ket{x} = \frac{1}{\sqrt{N}} \sum_{k=0}^{N-1}
e^{\frac{2\pi i x k}{N}} \ket{k}
$$
The analogy is clear: where the DFT maps a sequence of numbers to their
frequency components, the QFT maps a quantum state $\ket{x}$ to a
superposition of states $\ket{k}$, with amplitudes encoding the same
exponential structure. Applied to our periodic quantum state, the QFT
transforms the hidden period $r$ into measurable frequency information. The
QFT doesn't extract $r$ directly --- it extracts the frequency $1/r$, from
which we recover $r$ using the relation:
$$
\text{Period} = \frac{1}{\text{frequency}}
$$
A full treatment of the QFT would require its own article \cite{nielsen2010}.
What matters for our purposes is the following: after applying the QFT, the
amplitudes of our state are peaked near multiples of $j/r$, where $j$ is an
integer. This means that when we measure, we will most probably obtain a value
close to $j/r$, from which we can extract $r$ using the algorithm mentioned
in the next section.

\subsubsection{Continued Fraction Algorithm (CFA)}

After obtaining $j/r$, extracting $r$ from it is not straightforward. If
during your measurement, you measured that $j/r = 0.312$, extracting the
period from that is still a question that we need to answer.

For that, we use the \textbf{Continued Fraction Algorithm}. The best way to
explain the CFA is to give an example.

Let's assume that we found that $j/r = 0.312$.

We start by writing it as a fraction:
$$
0 + \frac{312}{1000}
$$

We then reverse the fraction and get this:
$$
= 0 + \frac{1}{\frac{1000}{312}} = 0 + \frac{1}{3 + \frac{8}{39}}
$$

And we repeat the same process until we get a $1$ in the numerator.
$$
0.312 = 0 + \frac{1}{3 + \frac{1}{4 + \frac{1}{1 + \frac{1}{7}}}}
$$

For each fraction, we get a better approximation for $0.312$:
$$
0.312 \approx \frac{1}{3}, \quad j = 1, \quad r = 3 \qquad \qquad 0.312 \approx \frac{1}{3 + \frac{1}{4}} = \frac{4}{13}, \quad j = 4, \quad r = 13
$$
% $$
% 0.312 \approx \frac{1}{3 + \frac{1}{4}} = \frac{4}{13}, \quad j = 4, \quad r = 13
% $$
$$
0.312 \approx \frac{1}{3 + \frac{1}{4 + \frac{1}{1}}} = \frac{5}{16},
\quad j = 5, \quad r = 16
$$

For a more detailed treatment of the algorithm and its post-quantum
implications, see \cite{watkins}.

We choose an approximation of $r$ such that it follows two conditions:

\begin{tcolorbox}[colframe=smjcrimson, colback=smjlightcrimson,
                  sharp corners, boxrule=0.5pt]
\begin{itemize}
  \item[1--] $r$ should be less than $N$
  \item[2--] $r$ should be an even number.
\end{itemize}
\end{tcolorbox}

If none of our approximations gives a valid $r$, we repeat the same process
from the beginning.

And with that, our quantum subroutine comes to an end. Let us take a step back
and appreciate what just happened. We started with a seemingly impossible
problem: factoring a large semiprime $N$. Rather than attacking it directly,
Shor's algorithm translated it into a period-finding problem --- one that a
quantum computer can solve efficiently. We made a random guess $g$, used the
quantum subroutine to find its period $r$, and from that period extracted the
prime factors of $N$ using nothing more than Euclid's algorithm and some
modular arithmetic.

Now that we have the full picture, we can revisit the main steps of the
algorithm and fill in the missing piece:

\begin{tcolorbox}[colframe=smjcrimson, colback=smjlightcrimson,
                  sharp corners, boxrule=0.5pt, breakable]
1 -- Make a guess, $g < N$ that shares no factors with $N$.\\[0.5em]
2 -- Use the quantum subroutine to find $r$ such that
     $g^r \equiv 1 \pmod{N}$:
\begin{itemize}
    \item Initialize two quantum registers: the first with $2n$ qubits in
          state $\ket{0}^{\otimes 2n}$, the second with $n$ qubits in
          state $\ket{1}$.
    \item Apply the Hadamard transform to the first register, creating an
          equal superposition of all states $\ket{x}$ for
          $x \in \{0, \ldots, 2^n - 1\}$.
    \item Apply $U$ to the full system:
          $U(\ket{x} \otimes \ket{y}) = \ket{x} \otimes \ket{y \cdot g^x \bmod N}$.
    \item Apply the QFT to the first register to extract frequency information.
    \item Measure the first register to obtain $j/r$, then apply the CFA
          to extract $r$.
\end{itemize}
\vspace{0.5em}
3 -- If $r$ is even, calculate $(g^{r/2} + 1)$ and $(g^{r/2} - 1)$.
     If $r$ is odd, go back to step 1.\\[0.5em]
4 -- Use Euclid's algorithm to find the greatest common divisor.
\end{tcolorbox}

This is the full picture of Shor's algorithm. But a natural question remains:
if this algorithm can break RSA, why is our data still safe today?
